%!TEX program = xelatex
%!TEX root = ../all.tex

\chapter{Clustering}

Clustering is one of the central tasks in unsupervised learning. Its defining characteristic is
that the goal of organising data into meaningful groups is pursued \emph{without} any
reliance on ground-truth labels. This stands in sharp contrast to classification, where the
training phase is guided by known targets and the learning objective is explicitly tied to
predictive accuracy. In clustering, the data arrive unlabeled and the aim is fundamentally
exploratory: we seek to uncover latent structure, recurring patterns, or meaningful groupings
that are only implicit in the geometry of the observations, with no external supervision to tell
us whether a proposed partition is correct or not. In practice, clustering serves as a
preliminary step for data understanding, for feature engineering, or as a component inside
larger pipelines---for instance, to compress data, to initialize mixture models, or to detect
anomalies.

\subsection*{Clustering Problem Setup}

Let $\mathbf{X}=(\mathbf{x}_1,\ldots,\mathbf{x}_n)$ be a dataset of $n$ points in
$\mathbb{R}^d$. A clustering task seeks to partition the dataset into $k$ clusters
$C_1,\ldots,C_k$ so that points falling in the same cluster are, in some sense,
\emph{similar}, while points belonging to different clusters are \emph{dissimilar}. The
precise meanings of ``similar'' and ``dissimilar'' are not universal: they depend on the
choice of a distance or similarity function and, more broadly, on the inductive bias we wish
to impose---compactness, separation, density, connectivity, and so on. For the methods
discussed in this chapter, similarity will be expressed primarily in terms of distances in
$\mathbb{R}^d$.

Rather than treating clustering as a purely combinatorial partitioning problem, it is often
more illuminating to view it as an \emph{optimization} problem: we define a quantitative
notion of cluster quality and then search for the grouping that optimises it. This perspective
makes the learning objective explicit and, as will become clear when we study $k$-means,
also suggests practical algorithmic strategies for finding good solutions.

\subsubsection*{Cluster representation}

A convenient way to represent a clustering is to make explicit both \emph{where} clusters
are located in feature space and \emph{which} points belong to each cluster. For this reason
we typically introduce two objects: a set of \textbf{prototypes} and an
\textbf{assignment matrix}.

Each cluster is summarised by a prototype $\mathbf{m}_j\in\mathbb{R}^d$, for
$j=1,\ldots,k$. Intuitively, $\mathbf{m}_j$ plays the role of a representative point for
cluster $C_j$; in many algorithms it corresponds to the ``center'' of the cluster, although
the precise meaning of ``center'' depends on the loss function and on the metric in use.

To encode membership, each data point $\mathbf{x}_i$ is associated with an assignment
vector $\mathbf{r}_i\in\{0,1\}^k$. The entry $r_{ij}$ indicates whether $\mathbf{x}_i$
belongs to cluster $C_j$:
\begin{myequation}
r_{ij}=
\begin{cases}
1 & \text{if } \mathbf{x}_i \text{ is assigned to cluster } C_j\\
0 & \text{otherwise.}
\end{cases}
\end{myequation}
This is a \alert{1-of-$k$ encoding}: for each $i$ exactly one component of $\mathbf{r}_i$
equals $1$ and all others equal $0$. In other words, the clustering is \emph{hard}: every
point is assigned to one and only one cluster. As we shall see later, probabilistic models
such as mixture models relax this constraint and allow \emph{soft} assignments, but
$k$-means is firmly in the hard-assignment setting.

\subsection*{Clustering approaches: partitional vs.\ hierarchical}

Although many clustering families exist, it is useful to distinguish two broad paradigms
that organise most classical methods: \emph{partitional} and \emph{hierarchical}
clustering. The difference is not only algorithmic; it also affects what kind of output the
method produces and what questions it is naturally suited to answer.

Partitional methods aim directly at producing a single, flat partition of the dataset into $k$
non-overlapping clusters. The number $k$ is typically fixed in advance (or selected by
model selection criteria), and the algorithm searches for a clustering that optimises a cost
function $J$ measuring cluster quality. Under this view, clustering closely resembles
parameter estimation: we are simultaneously choosing the prototypes---or other parameters
describing clusters---and the assignments of points to clusters. A key advantage of
partitional methods is computational scalability: with appropriate implementations, they are
often effective on large datasets. Typical representatives include $k$-means, $k$-medoids,
and, with a probabilistic interpretation, mixture models.

Hierarchical methods do not return a single flat partition. Instead, they produce a nested
family of partitions organised as a tree, commonly visualised as a \emph{dendrogram}.
Moving along the tree corresponds to changing the granularity of the description: at the
bottom we find many small clusters, and as we move upward, clusters merge in the
agglomerative approach (or, dually, large clusters split into smaller ones in the divisive
approach). One practical benefit is that we do not need to commit to a specific $k$ at the
outset; any desired number of clusters can be obtained by ``cutting'' the dendrogram at a
chosen height. This richer description is often informative, but it typically comes with
higher computational cost. Moreover, hierarchical decisions made early in the process
cannot be revised later, which makes the method sensitive to the choice of linkage criterion
and distance metric.

These two paradigms are not in competition: they are complementary, and the appropriate
choice depends on the goals of the analysis. When speed and scalability matter, or when
there is a reasonable prior belief about the number of clusters, partitional methods are
often preferred. When the primary goal is to understand the multi-scale structure of the
data, or when $k$ is genuinely unknown, hierarchical methods offer a richer picture.

\subsection*{The $k$-Means Algorithm}

Among partitional clustering methods, $k$-means occupies a special place because of its
simplicity, speed, and clarity of interpretation. Despite being conceptually straightforward,
it embodies many themes that recur in more advanced unsupervised learning algorithms:
optimisation under discrete assignments, alternating minimisation, convergence to local
optima, and sensitivity to initialization. Understanding $k$-means in depth therefore
provides a strong foundation for the more sophisticated methods discussed later in this
chapter.

Given a dataset $\mathbf{X}$ and a target number of clusters $k$, $k$-means searches
simultaneously for an assignment matrix $R=(r_{ij})$ and a set of prototypes
$M=(\mathbf{m}_1,\ldots,\mathbf{m}_k)$ that minimise the following objective:
\begin{myequation}
J(R,M)=\sum_{i=1}^{n}\sum_{j=1}^{k} r_{ij}\|\mathbf{x}_i-\mathbf{m}_j\|^2.
\end{myequation}
This quantity is known as the \alert{within-cluster sum of squares}. It measures how tightly
each cluster concentrates around its prototype: because $r_{ij}$ is zero for all clusters
except the one to which $\mathbf{x}_i$ is assigned, each point contributes only through its
squared distance to the prototype of its own cluster.

\paragraph{Interpretation.}
The objective makes the inductive bias of $k$-means explicit. Minimising $J$ encourages
\emph{compact} clusters in Euclidean space: points should lie close to their cluster center,
and a good clustering is one in which most of the squared distances $\|\mathbf{x}_i-
\mathbf{m}_j\|^2$ (for the assigned $j$) are small. This bias fits well when clusters are
roughly spherical---or, more precisely, when cluster structure is well captured by Euclidean
distance and means---but it becomes restrictive in the presence of elongated, curved, or
density-based cluster shapes. Keeping this limitation in mind is important when deciding
whether $k$-means is the right tool for a given problem.

\subsubsection*{Algorithm: alternating optimisation}

Jointly optimising $R$ and $M$ is difficult because $R$ is discrete and the objective is not
convex in the pair $(R,M)$. The classical $k$-means algorithm addresses this by
\alert{alternating optimisation}: it iteratively improves one set of variables while keeping
the other fixed. Each individual step is simple and has a closed-form solution, but together
they produce a monotone decrease of the objective that guarantees convergence.

The assignment step minimises $J$ with respect to $R$, keeping $M$ fixed. With the
prototypes fixed, the best assignment for each point is obtained by choosing the nearest
prototype in squared Euclidean distance:
\begin{myequation}
r_{ij}=
\begin{cases}
1 & \text{if } j=\argmin{\ell}\ \|\mathbf{x}_i-\mathbf{m}_{\ell}\|^2\\
0 & \text{otherwise.}
\end{cases}
\end{myequation}
Geometrically, this step partitions the space into Voronoi-like regions around the current
prototypes: each point is assigned to the region of the closest prototype, and the decision
boundaries are piecewise linear hyperplanes equidistant from neighbouring prototypes.

The update step then minimises $J$ with respect to $M$, keeping $R$ fixed. For a given
cluster index $j$, differentiating $J$ with respect to $\mathbf{m}_j$ and setting the gradient
to zero yields a straightforward calculation:
\begin{myalign}
\frac{\partial J}{\partial \mathbf{m}_j} &= 2\sum_{i=1}^{n} r_{ij}(\mathbf{m}_j-\mathbf{x}_i)=0\\
\Rightarrow \mathbf{m}_j &= \frac{\sum_{i=1}^{n} r_{ij}\mathbf{x}_i}{\sum_{i=1}^{n} r_{ij}}
= \frac{1}{n_j}\sum_{i\in C_j}\mathbf{x}_i,
\end{myalign}
where $n_j=\sum_{i=1}^{n} r_{ij}$ is the number of points assigned to cluster $C_j$. The
optimal prototype is therefore the \emph{mean} of all points currently assigned to that
cluster---which is precisely why the algorithm is called $k$-\textit{means}. Geometrically,
the centroid is the unique point that minimises the sum of squared distances to any given
set of points, a classical result from elementary statistics.

\subsubsection*{Convergence properties}

A useful feature of $k$-means is that it is guaranteed to make progress in terms of the
objective at every iteration. After each full cycle of assignment followed by update, the cost
cannot increase:
\begin{myequation}
J^{(t+1)}\leq J^{(t)}.
\end{myequation}
The argument is embedded directly in the construction of the two steps. The assignment
step chooses, for each point, the cluster that minimises its contribution to $J$ given the
current prototypes: any reassignment of a point to a closer prototype cannot increase the
total cost. The update step then replaces each prototype by the centroid of its currently
assigned cluster: since the centroid minimises the within-cluster sum of squares, updating
the prototypes can only reduce $J$ further. Because $J(R,M)\geq 0$ always holds, this
monotone decrease implies that the algorithm converges in a finite number of steps, since
there are only finitely many possible assignment matrices and the cost strictly decreases
unless a fixed point has already been reached.

It is important, however, not to overstate this result. The objective is not convex, and
monotone descent does \emph{not} imply convergence to the global optimum. In practice,
$k$-means typically converges to a \textbf{local} minimum of $J$, and different
initializations can lead to qualitatively different solutions. This is not a pathological edge
case: it is a fundamental aspect of the optimisation landscape induced by discrete
assignments and squared distances. In datasets with complex structure or with clusters of
comparable within-cluster variance, the algorithm may converge to partitions that differ
substantially from the globally optimal one.

\subsubsection*{Initialization and stopping criteria}

The simplest initialization strategy selects $k$ prototypes uniformly at random among the
data points, and then begins alternating between assignment and update. Because different
random starts can lead to different local minima, it is common practice to run the algorithm
multiple times from different starting configurations and retain the final clustering with the
smallest attained cost. More sophisticated initialization schemes, such as $k$-means++,
improve over pure random initialization by choosing starting prototypes that are spread out
across the dataset, which tends to produce better solutions and faster convergence.

Stopping can be decided in several equivalent ways. One may halt after a fixed number of
iterations, or when the cost improvement becomes negligible, i.e., when
$|J^{(t)}-J^{(t-1)}|<\epsilon$ for a chosen tolerance $\epsilon$. Another natural stopping
criterion is to stop when the assignments no longer change between iterations: this indicates
that the algorithm has reached a fixed point at which the Voronoi partition induced by the
current prototypes is consistent with those same prototypes being the centroids of their
respective clusters.

\subsection*{$k$-Means Example}

The following sequence of figures illustrates the evolution of $k$-means on a
two-dimensional dataset with $k=3$. The geometric meaning of the alternating steps
becomes visible: assignments carve the plane into decision regions, while prototype updates
move the centers toward the centroids of the currently assigned points. Even in this simple
example, one can observe how the algorithm quickly reorganises the prototypes during the
first few iterations and then settles into a stable configuration.

\begin{figure}[h]
\centering
\begin{subfigure}{0.32\textwidth}
\centering
\includegraphics[width=\textwidth]{kmeans0.pdf}
\caption{Original dataset}
\end{subfigure}
\hfill
\begin{subfigure}{0.32\textwidth}
\centering
\includegraphics[width=\textwidth]{kmeans1.pdf}
\caption{Random initialization}
\end{subfigure}
\hfill
\begin{subfigure}{0.32\textwidth}
\centering
\includegraphics[width=\textwidth]{kmeans2.pdf}
\caption{After iteration 1}
\end{subfigure}
\end{figure}

\begin{figure}[h]
\centering
\begin{subfigure}{0.32\textwidth}
\centering
\includegraphics[width=\textwidth]{kmeans3.pdf}
\caption{After iteration 2}
\end{subfigure}
\hfill
\begin{subfigure}{0.32\textwidth}
\centering
\includegraphics[width=\textwidth]{kmeans4.pdf}
\caption{After iteration 3}
\end{subfigure}
\hfill
\begin{subfigure}{0.32\textwidth}
\centering
\includegraphics[width=\textwidth]{kmeans5.pdf}
\caption{Final clustering}
\end{subfigure}
\end{figure}

\medskip
\noindent Figure caption: The algorithm alternates between assigning points to the nearest
prototype---thus inducing decision regions in the plane---and recomputing prototypes as
cluster means. In many practical datasets, convergence is reached in a small number of
iterations.

\subsubsection*{Cost curve}

A complementary perspective on convergence is obtained by monitoring the objective value
over iterations. The monotone decrease predicted by theory typically appears clearly in
practice, with a rapid initial drop followed by progressively smaller improvements as the
algorithm approaches a stable configuration. Figure~\ref{f:cf} shows a representative cost
curve: the steep descent in the first few iterations reflects large reorganisations of the
assignment matrix, while the subsequent plateau indicates that only marginal adjustments
are being made and the algorithm is close to a fixed point.

\begin{figure}[h]
\centering
\includegraphics[width=0.6\textwidth]{kmeansobj.pdf}
\caption{Cost function $J(R,M)$ over iterations. The cost decreases monotonically and
stabilises as the algorithm converges.}\label{f:cf}
\end{figure}

\subsection*{Choosing the number of clusters $k$}

The value of $k$ is a crucial hyperparameter in partitional clustering: it controls the
granularity of the partition and strongly affects the interpretation of the result. If $k$ is too
small, distinct groups may be merged and the clustering will be overly coarse. If $k$ is too
large, the algorithm may fragment meaningful clusters into smaller pieces, often fitting
noise rather than genuine structure. Selecting $k$ is therefore not a minor technical detail
but a substantive modelling decision.

\subsubsection*{Elbow method}

A simple and widely used heuristic is the \alert{elbow method}. The idea is to run
$k$-means for a range of values of $k$---say, $k=1,2,\ldots,10$---and to plot the achieved
cost $J(k)$ as a function of $k$. Because adding more clusters increases the flexibility of
the model, the cost necessarily decreases as $k$ grows. However, there typically exists a
point beyond which the improvement becomes marginal: the ``elbow'' of the curve, where
the decrease starts to flatten, is taken as a reasonable compromise between fit quality and
model complexity.

\begin{figure}[h]
\centering
\includegraphics[width=0.6\textwidth]{clusteringgraph.pdf}
\caption{Elbow method: cost vs.\ number of clusters. The elbow indicates a good choice of
$k$ where additional clusters provide diminishing returns.}
\end{figure}

While intuitive, the elbow method is not always decisive: in some datasets the curve
decreases smoothly without a clear kink, making the ``elbow'' difficult to identify
unambiguously. In those cases, it is useful to complement this heuristic with more
principled, model-based criteria, which we now discuss.

\subsubsection*{Other model selection criteria}

Information-theoretic criteria provide a systematic way to balance fit quality against
model complexity, making explicit the trade-off between explaining the data well and
avoiding unnecessary parameters. To apply such criteria, we need to embed $k$-means in
a probabilistic framework.

Although $k$-means is not inherently probabilistic, it can be embedded in a simple
generative model. Assume that, conditional on belonging to cluster $j$, points are
distributed according to an isotropic Gaussian centered at $\mathbf{m}_j$ with covariance
$\sigma^2\mathbf{I}$:
\begin{myequation}
p(\mathbf{x}_i \mid \text{cluster } j)=\frac{1}{(\sqrt{2\pi}\sigma)^d}
\exp\!\left(-\frac{\|\mathbf{x}_i-\mathbf{m}_j\|^2}{2\sigma^2}\right).
\end{myequation}
Given hard assignments, the likelihood of the full dataset is
\begin{myequation}
L=\prod_{i=1}^{n} p(\mathbf{x}_i \mid C_{j_i}),
\end{myequation}
where $j_i$ denotes the cluster assigned to $\mathbf{x}_i$. Taking logarithms gives
\begin{myequation}
\ln L=\sum_{i=1}^{n}\left[-d\ln(\sqrt{2\pi}\sigma)
-\frac{\|\mathbf{x}_i-\mathbf{m}_{j_i}\|^2}{2\sigma^2}\right].
\end{myequation}
Ignoring constants that do not depend on $k$, we obtain
\begin{myequation}
\ln L \propto -\frac{1}{2\sigma^2}\sum_{i=1}^{n}\|\mathbf{x}_i-\mathbf{m}_{j_i}\|^2
= -\frac{J(R,M)}{2\sigma^2}.
\end{myequation}
This reveals a key connection: minimising the $k$-means objective $J(R,M)$ is equivalent
to maximising the likelihood under the isotropic Gaussian assumption. This equivalence is
not merely a mathematical curiosity; it provides the justification for applying standard
model selection tools to the $k$-means problem.

\paragraph{Akaike Information Criterion (AIC).}
AIC trades off goodness of fit against complexity through the formula
\begin{myequation}
\text{AIC}(k)=2\cdot(\text{number of parameters})-2\ln L.
\end{myequation}
In $k$-means with $k$ prototypes in $d$ dimensions, the number of mean parameters is
$kd$, so
\begin{myequation}
\text{AIC}(k)=2kd-2\ln L=2kd+\frac{J(R,M)}{\sigma^2}.
\end{myequation}
The interpretation is transparent: increasing $k$ typically decreases $J(R,M)$ (better fit),
but simultaneously increases the penalty $2kd$ (greater model complexity). AIC balances
these two competing effects, and the selected $k$ is the one that \emph{minimises} it.

\paragraph{Bayesian Information Criterion (BIC).}
BIC is motivated by Bayesian arguments and imposes a stronger penalty on model
complexity:
\begin{myequation}
\text{BIC}(k)=(\text{number of parameters})\cdot \ln n - 2\ln L.
\end{myequation}
For $k$-means:
\begin{myequation}
\text{BIC}(k)=kd\cdot \ln n+\frac{J(R,M)}{\sigma^2}.
\end{myequation}
Compared to AIC, the penalty term $kd\ln n$ grows with the sample size $n$, meaning that
with more data BIC becomes increasingly conservative and tends to prefer simpler models
unless the fit improvement is clearly justified by the evidence. In practice, when
$n$ is large, BIC often selects fewer clusters than AIC, reflecting its stronger guard against
overfitting.

\paragraph{Silhouette coefficient.}
The silhouette coefficient is a distance-based criterion that does not require a probabilistic
model. It evaluates clustering quality by combining two complementary intuitions: points
should be close to other points in their own cluster (compactness) and far from points in
other clusters (separation). Consider a point $\mathbf{x}_i$ assigned to cluster $C_j$, and
define $a_i$ as the average distance from $\mathbf{x}_i$ to all other points in the same
cluster:
\begin{myequation}
a_i=\frac{1}{|C_j|-1}\sum_{\mathbf{x}_\ell\in C_j,\ \ell\neq i} d(\mathbf{x}_i,\mathbf{x}_\ell).
\end{myequation}
A small $a_i$ indicates that $\mathbf{x}_i$ sits well within its cluster. Next, define $b_i$
as the smallest average distance from $\mathbf{x}_i$ to the points of any other cluster:
\begin{myequation}
b_i=\min_{k\neq j}\ \frac{1}{|C_k|}\sum_{\mathbf{x}_\ell\in C_k} d(\mathbf{x}_i,\mathbf{x}_\ell).
\end{myequation}
A large $b_i$ means that the nearest alternative cluster is still far away, which is
desirable. The silhouette value for $\mathbf{x}_i$ combines these two quantities as follows:
\begin{myequation}
s_i=\frac{b_i-a_i}{\max(a_i,b_i)}.
\end{myequation}
By construction, $s_i\in[-1,1]$. Values close to $+1$ indicate that the point is well
clustered (small $a_i$ and large $b_i$); values close to $0$ indicate a borderline point
that lies between two clusters; and negative values suggest that $\mathbf{x}_i$ may have
been assigned to the wrong cluster. The global quality of a clustering is then summarised by
the average silhouette coefficient:
\begin{myequation}
S=\frac{1}{n}\sum_{i=1}^{n} s_i,
\end{myequation}
and a common strategy is to select the value of $k$ that maximises $S$. Unlike AIC and
BIC, the silhouette does not assume any generative model; it is therefore applicable
whenever a distance function between points is available, even when no probabilistic
interpretation is attached to the clusters.

Beyond these quantitative criteria, \textbf{domain knowledge} remains one of the most
reliable guides: if the application provides strong prior expectations about how many groups
should exist or what they should represent, that information can and should inform the
choice of $k$. Quantitative criteria and domain expertise are best used in conjunction rather
than as substitutes for one another.

\subsection*{Advantages and limitations of $k$-means}

It is useful to conclude this section with a balanced assessment of when $k$-means is an
appropriate choice and when its inherent assumptions become restrictive.

On the side of advantages, $k$-means is simple to implement and computationally
efficient. With a per-iteration cost that scales as $O(nkd)$ under standard assumptions, it
remains practical even for large datasets and serves as a common baseline and a frequent
component in larger systems. It performs well when clusters are compact, convex, and
roughly spherical in Euclidean space, which is a reasonable description of many real-world
datasets.

On the side of limitations, the same design choices that make $k$-means fast also define
its boundaries. Because it optimises a non-convex objective, it is sensitive to initialisation
and may converge to poor local minima---a risk that grows with the number of clusters and
the complexity of the data. It requires specifying $k$ in advance, which is not always
natural in truly exploratory settings where the number of groups is itself a quantity of
interest. Most importantly, the squared Euclidean objective implicitly assumes clusters that
are similar in size and shape; it struggles with non-convex, elongated, or manifold-like
clusters. Furthermore, because prototypes are means, outliers can exert a disproportionate
influence and distort cluster centers, leading to unstable or misleading partitions.

These limitations motivate the methods introduced in the rest of this chapter:
hierarchical clustering, which avoids fixing $k$ in advance and produces a richer
multi-scale description of the data, and probabilistic mixture models, which relax the
spherical-cluster assumption and replace deterministic assignments with a principled
treatment of uncertainty.

\subsection*{Hierarchical Clustering}

\subsubsection*{Overview}

In many exploratory scenarios, the question of interest is not only \emph{``how do I
partition the data into $k$ groups?''} but also \emph{``which groupings are plausible at
different levels of detail?''} Hierarchical clustering addresses precisely this point. Rather
than committing to a single flat partition, it produces a \alert{tree-structured} description
of the dataset that encodes a whole \emph{family} of partitions, ranging from very
fine-grained (many small clusters) to very coarse (a single large cluster). This output is
visualised as a \alert{dendrogram}, a diagram in which every cut at a fixed height
corresponds to a distinct partition. The conceptual difference with respect to $k$-means is
therefore fundamental: instead of requiring the analyst to choose a granularity in advance,
hierarchical clustering defers that decision and presents the full spectrum of plausible
groupings together.

The hierarchy can be constructed in two symmetric ways. In the \textbf{agglomerative}
(bottom-up) approach, we begin with the most fragmented description possible---each
data point forming its own cluster---and then progressively merge clusters according to a
chosen notion of similarity. In the \textbf{divisive} (top-down) approach, we start from the
opposite extreme---all points in a single cluster---and repeatedly split clusters until every
point is in its own singleton. Although both perspectives are meaningful, agglomerative
hierarchical clustering is more commonly used in practice, largely because it admits simpler
and more efficient implementations. Divisive procedures, by contrast, often require solving
a harder optimisation problem at each split, which makes them computationally demanding.

To make these ideas precise, let $\mathbf{X}=(\mathbf{x}_1,\ldots,\mathbf{x}_n)$ with
$\mathbf{x}_i\in\mathbb{R}^d$, and let $d(\cdot,\cdot)$ denote a distance (or
dissimilarity) between individual points. Since hierarchical clustering builds clusters by
comparing not only single points but also sets of points, a key modelling choice is how to
lift the point-to-point distance $d$ into a distance between clusters. This is the role of the
\emph{linkage criterion}, which is introduced in the next section.

\subsubsection*{Agglomerative hierarchical clustering}

Agglomerative clustering begins from $n$ singleton clusters and merges them step by step
until only a single cluster remains. The procedure is greedy: at each iteration it merges the
two clusters that are currently the most similar---equivalently, the least distant---according
to the chosen linkage criterion. Formally, let
$\mathcal{C}^{(t)}=\{C_1^{(t)},\ldots,C_{m_t}^{(t)}\}$ be the current collection of
clusters at iteration $t$, where $m_t$ decreases from $n$ down to $1$.

The algorithm is initialised by setting $\mathcal{C}^{(0)}=\big\{\{\mathbf{x}_1\},\ldots,
\{\mathbf{x}_n\}\big\}$, i.e., every point is its own cluster. At each subsequent iteration
$t=0,1,\ldots,n-2$, the inter-cluster distances $d_{\text{link}}(C_a^{(t)},C_b^{(t)})$ are
computed for all pairs of distinct clusters in $\mathcal{C}^{(t)}$, the closest pair is
identified,
\begin{myequation}
(C_p^{(t)},C_q^{(t)}) \in \argmin{C_a^{(t)}\neq C_b^{(t)}} d_{\text{link}}(C_a^{(t)},C_b^{(t)}),
\end{myequation}
and these two clusters are merged into a new cluster
$C_{\text{new}}^{(t+1)}=C_p^{(t)}\cup C_q^{(t)}$, replacing them in the collection. The
full sequence of merges, together with the distance at which each merge occurred,
constitutes the dendrogram.

This greedy construction has an important consequence that is worth emphasising: once
two clusters have been merged, the decision is never revisited. There is no mechanism for
correcting a premature merge. For this reason, the choice of linkage criterion is not a
minor technicality but the main mechanism controlling the geometry of the resulting
hierarchy, and different choices can lead to qualitatively different dendrograms even when
the underlying point-to-point distance is the same.

\subsubsection*{Cluster similarity and linkage criteria}

A linkage criterion specifies how to compute a distance between two clusters from the
underlying distance between individual points. Different criteria correspond to different
implicit notions of what a ``good'' cluster should look like. In all cases, $C_1$ and $C_2$
denote two clusters (sets of points) and $d(\mathbf{x}_i,\mathbf{x}_j)$ is the base distance
between points, most often the Euclidean distance.

\paragraph{Single linkage.}
Single linkage defines the distance between two clusters as the distance between their
closest pair of points:
\begin{myequation}
d_{\text{SL}}(C_1,C_2)=\min_{\mathbf{x}_i\in C_1,\ \mathbf{x}_j\in C_2} d(\mathbf{x}_i,\mathbf{x}_j).
\end{myequation}
This criterion treats two clusters as similar as soon as \emph{any} bridge of nearby points
exists between them. As a result, it often produces long, chain-like clusters, a
phenomenon known as the \emph{chaining effect}. This behaviour can be desirable when
the true clusters are non-convex and connected by paths of close neighbours, but it also
makes single linkage sensitive to noise and outliers: a single stray point can connect two
otherwise well-separated groups, causing them to merge prematurely.

\paragraph{Complete linkage.}
Complete linkage takes the opposite perspective and measures the distance between two
clusters through their farthest pair of points:
\begin{myequation}
d_{\text{CL}}(C_1,C_2)=\max_{\mathbf{x}_i\in C_1,\ \mathbf{x}_j\in C_2} d(\mathbf{x}_i,\mathbf{x}_j).
\end{myequation}
Under this criterion, two clusters are considered similar only if \emph{all} points in one
cluster are relatively close to \emph{all} points in the other. This tends to create compact,
relatively spherical clusters and largely eliminates the chaining effect. The drawback is
that the maximum distance is determined by the farthest pair of points, which may be
atypical, so complete linkage can be overly conservative when clusters are irregular in
shape.

\paragraph{Group average linkage.}
Group average linkage (also known as average linkage) balances the previous two
extremes by using the average pairwise distance across all pairs of points from different
clusters:
\begin{myequation}
d_{\text{GA}}(C_1,C_2)=\frac{1}{|C_1|\cdot |C_2|}
\sum_{\mathbf{x}_i\in C_1}\sum_{\mathbf{x}_j\in C_2} d(\mathbf{x}_i,\mathbf{x}_j).
\end{myequation}
By averaging over all cross-cluster pairs, this criterion is often more robust than single
linkage and less stringent than complete linkage. It merges clusters when they are close
\emph{on average}, yielding a natural compromise between connectivity and compactness.
In practice, average linkage is frequently the preferred choice when there is no strong
a priori reason to favour either extreme.

The linkage choice can dramatically change the resulting dendrogram even when the
underlying point distance $d$ is held fixed. Hierarchical clustering is therefore best
understood as a \emph{framework} rather than a single algorithm: the specific hierarchy
produced depends jointly on the point distance and the linkage criterion, each encoding a
distinct modelling assumption about the shape and structure of clusters.

\subsubsection*{Interpreting a dendrogram}

A dendrogram is a rooted binary tree that summarises the entire merging process. Its
leaves correspond to individual data points, each internal node corresponds to a merge
between two clusters, and the height (vertical position) of each internal node encodes the
linkage distance $d_{\text{link}}$ at which the corresponding merge took place. The height
therefore has an immediate and interpretable meaning: merges occurring at low heights
involve very similar clusters, while merges occurring at high heights indicate that the
algorithm had to join rather dissimilar groups because all more natural alternatives had
already been exhausted. In this sense, the dendrogram provides a visual map of similarity
relationships across the entire dataset and can be read as a multi-scale description of its
structure.

A practical advantage of this representation is that it naturally induces a partition for
\emph{any} desired number of clusters without requiring a new run of the algorithm.
Cutting the dendrogram horizontally at height $h$ produces as many clusters as there are
connected subtrees below the cut: a lower cut yields more clusters (a finer partition),
while a higher cut yields fewer (a coarser partition). This flexibility allows the analyst to
explore plausible values of $k$ a posteriori by inspecting the dendrogram, looking for
heights at which large gaps appear---indicating that a merge at that level combines two
genuinely dissimilar groups---and deciding on a cut accordingly.

\subsubsection*{Example dendrogram}

\begin{figure}[h]
\centering
\includegraphics[width=0.7\textwidth]{iris2.pdf}
\caption{Dendrogram from hierarchical clustering on a dataset. The height of each merge
indicates the dissimilarity between merged clusters. Cutting at different heights yields
different numbers of clusters.}
\end{figure}

\subsubsection*{Properties and practical considerations}

Hierarchical clustering has several characteristic properties worth summarising. First, it
does not require specifying $k$ in advance: the dendrogram implicitly contains solutions for
all possible values of $k$, which is particularly useful in exploratory analysis where the
number of groups is not known a priori. Second, in the agglomerative setting, the sequence
of merge heights is monotone non-decreasing: the algorithm always merges the closest
available pair, so as clusters grow the next merge can only occur at the same or higher
dissimilarity. This monotonicity is precisely what makes the dendrogram a consistent and
well-defined representation of the process.

At the same time, hierarchical clustering is sensitive to the linkage criterion. Because
merges are irreversible, two different linkage choices can produce qualitatively different
hierarchies on the same data. When interpretability is the primary goal, it is therefore often
informative to compare dendrograms obtained with different linkages and assess how robust
the observed structure is to this modelling choice.

Regarding the practical trade-off with $k$-means: when the dataset is large and a specific
number of clusters is known or can be reliably estimated, $k$-means is often preferred
because of its speed and scalability. Hierarchical clustering, by contrast, is especially
valuable when the primary goal is to \textbf{understand structure}: it supports
small-to-moderate datasets, helps diagnose whether multiple granularities are plausible,
and offers an interpretable, parameter-free picture of how clusters relate to each other
across scales.

\section*{Probabilistic Clustering: Mixture Models}

The clustering methods introduced so far---$k$-means and hierarchical clustering---are
primarily \emph{geometric}: they define clusters through distances and partitions of the
feature space, but they make no explicit statement about \emph{how the data might have
been generated}. In many applications, however, it is natural and productive to assume that
the observed points are realisations of an underlying random process. From this
perspective, a cluster is not merely a region of space but rather a \emph{population}
described by a probability distribution. This motivates a \alert{probabilistic approach} to
clustering, in which we posit a statistical model for the data and interpret clusters as the
latent components of that model.

\subsection*{The mixture model framework}

\subsubsection*{Motivation}

The central idea is to describe the data distribution as a \alert{mixture} of simpler
component distributions. Generating a point $\mathbf{x}$ from a mixture can be thought of
as a two-stage process: first choose a component (which we interpret as a cluster)
according to some prior probabilities, then sample the point from the distribution associated
with the chosen component. This viewpoint introduces two important differences with
respect to $k$-means. First, cluster membership becomes a \emph{latent} random variable
rather than a deterministic label: we cannot observe which component generated a given
point, and the goal of learning is to recover this latent structure from the data. Second,
because the model is probabilistic, we can quantify uncertainty: a point may plausibly
belong to more than one cluster with non-negligible probability, and the model makes this
ambiguity explicit rather than suppressing it.

This probabilistic formulation brings several benefits beyond producing a partition of the
dataset. It naturally yields \emph{soft assignments}, supports likelihood-based reasoning
and model comparison, and enables principled criteria such as AIC and BIC to guide the
choice of the number of components. Moreover, because we can select different distribution
families for the components, the same conceptual framework extends seamlessly from
continuous measurements to count data, binary vectors, and many other data types.

\subsubsection*{Mathematical definition}

A \alert{mixture model} is defined as a weighted sum of component distributions:
\begin{myequation}
p(\mathbf{x}\mid\boldsymbol{\pi},\boldsymbol{\Theta})=\sum_{k=1}^{K}\pi_k\,q(\mathbf{x}\mid\boldsymbol{\theta}_k),
\end{myequation}
where $K$ is the number of components, $\pi_k$ are the \alert{mixing coefficients}
satisfying $\sum_{k=1}^{K}\pi_k=1$ and $0\leq \pi_k\leq 1$, and
$q(\mathbf{x}\mid\boldsymbol{\theta}_k)$ is the $k$-th component distribution with
parameters $\boldsymbol{\theta}_k$. We collect the mixing coefficients in
$\boldsymbol{\pi}=(\pi_1,\ldots,\pi_K)$ and the component parameters in
$\boldsymbol{\Theta}=(\boldsymbol{\theta}_1,\ldots,\boldsymbol{\theta}_K)$.

The interpretation is straightforward but important. The coefficient $\pi_k$ represents the
probability that a randomly sampled point was generated by component $k$
\emph{before} observing the point itself: it is the \emph{prior} prevalence of cluster $k$
in the population. The function $q(\cdot\mid\boldsymbol{\theta}_k)$ then describes how
points are distributed \emph{within} that cluster. Because the mixture is a convex
combination of valid densities or mass functions, it is itself a valid probability distribution:
it is nonnegative and integrates (or sums) to $1$ over its domain.

Although many expositions restrict attention to mixtures whose components all belong to
the same parametric family, the definition does not require this. Components from different
families can in principle be combined. In practice, however, using a single family is both
simpler and better motivated: it allows clusters to differ only through their parameters
(e.g., their mean and covariance in the Gaussian case), making the estimated components
directly comparable.

\subsubsection*{Choice of component distributions}

The component family should be chosen to match the nature of the observations. For
continuous real-valued data in $\mathbb{R}^d$, the canonical choice is the multivariate
Gaussian, leading to Gaussian mixture models. For count data, Poisson mixtures are
natural; for binary vectors, mixtures of Bernoulli distributions provide a basic but flexible
model. Using the same family for all components simplifies parameter estimation, because
each component can be updated using the same type of sufficient statistics, while still
allowing considerable flexibility through different parameter values. One should keep in
mind, however, that the mixture family is a modelling assumption: if the chosen
components poorly match the true data-generating mechanism, the inferred clusters may be
unstable or difficult to interpret.

\subsection*{Clustering via mixture models}

The connection between mixture models and clustering becomes explicit once we introduce
a latent variable $z\in\{1,\ldots,K\}$ indicating which component generated a given
observation $\mathbf{x}$. Under the generative story, one first draws
$z\sim\text{Categorical}(\boldsymbol{\pi})$, then draws
$\mathbf{x}\sim q(\cdot\mid\boldsymbol{\theta}_z)$. Marginalising out the unobserved $z$
recovers the mixture formula introduced above. In this setting, clustering amounts to
inferring---at least approximately---which component generated each observation, given
the fitted parameters. Because $z$ is latent, we do not obtain a single deterministic
assignment unless we explicitly force one; instead, we obtain \emph{posterior probabilities}
of membership, which carry information about ambiguity that a hard assignment would
discard.

\subsubsection*{Soft vs.\ hard assignments and responsibilities}

For a given point $\mathbf{x}$, Bayes' theorem gives the posterior probability that it was
generated by component $k$:
\begin{myequation}
\gamma_k(\mathbf{x})=p(z=k\mid\mathbf{x})=
\frac{\pi_k\,q(\mathbf{x}\mid\boldsymbol{\theta}_k)}
{\sum_{j=1}^{K}\pi_j\,q(\mathbf{x}\mid\boldsymbol{\theta}_j)}.
\end{myequation}
The quantity $\gamma_k(\mathbf{x})$ is called the \alert{responsibility} of component $k$ for
$\mathbf{x}$. It satisfies $0\leq \gamma_k(\mathbf{x})\leq 1$ and
$\sum_{k=1}^{K}\gamma_k(\mathbf{x})=1$, so it defines a proper probability distribution
over clusters for the specific point $\mathbf{x}$. If, for instance,
$\gamma_1(\mathbf{x})=0.9$ and $\gamma_2(\mathbf{x})=0.1$, then $\mathbf{x}$ is
strongly associated with component $1$, yet the model acknowledges a non-negligible
probability of belonging to component $2$. This is precisely the kind of uncertainty
information that is absent in hard clustering methods like $k$-means.

When a hard partition is needed, one can still obtain it from a mixture model by assigning
each point to the component with the largest responsibility---a maximum a posteriori
decision rule. However, the soft assignments are often more informative, especially near
overlapping regions where the data do not clearly support a crisp decision boundary.

\subsection*{Parameter estimation by maximum likelihood}

\subsubsection*{The likelihood and the missing-data difficulty}

Given a dataset $\mathbf{X}=(\mathbf{x}_1,\ldots,\mathbf{x}_n)$, the likelihood under the
mixture model is
\begin{myequation}
L(\boldsymbol{\pi},\boldsymbol{\Theta}\mid\mathbf{X})
=\prod_{i=1}^{n}p(\mathbf{x}_i\mid\boldsymbol{\pi},\boldsymbol{\Theta})
=\prod_{i=1}^{n}\sum_{k=1}^{K}\pi_k\,q(\mathbf{x}_i\mid\boldsymbol{\theta}_k),
\end{myequation}
and the corresponding log-likelihood is
\begin{myequation}
\ell(\boldsymbol{\pi},\boldsymbol{\Theta}\mid\mathbf{X})
=\sum_{i=1}^{n}\log\!\left(\sum_{k=1}^{K}\pi_k\,q(\mathbf{x}_i\mid\boldsymbol{\theta}_k)\right).
\end{myequation}
Maximising $\ell$ with respect to $\boldsymbol{\pi}$ and $\boldsymbol{\Theta}$ is, in
general, nontrivial. The logarithm contains a sum over components, which couples the
parameters in a nonlinear way, and the mixing coefficients must satisfy the simplex
constraint $\sum_{k=1}^{K}\pi_k=1$. Setting the gradient to zero does not yield a closed-
form solution. The conceptual root of this difficulty is that the data are \alert{incomplete}:
we observe each $\mathbf{x}_i$, but we do not observe which component $z_i$ generated
it. If the component indicators $(z_1,\ldots,z_n)$ were known, parameter estimation would
reduce to standard maximum likelihood for each component separately, plus a simple
frequency estimate of the mixing proportions. The EM algorithm exploits precisely this
intuition, iteratively constructing a tractable surrogate for the incomplete-data likelihood.

\subsection*{The expectation--maximization (EM) algorithm for mixtures}

EM is a general strategy for maximum likelihood estimation in latent-variable models that
will be introduced in its full generality later. For mixture models, it suffices to observe that
it alternates between two steps: computing posterior responsibilities (a probabilistic
``completion'' of the missing labels) and updating parameters as if those soft labels were
observed.

\subsubsection*{Computing responsibilities (the E step)}

Given current parameter estimates $\boldsymbol{\pi}^{(t)}$ and $\boldsymbol{\Theta}^{(t)}$,
the E step computes, for each point $\mathbf{x}_i$ and each component $k$,
\begin{myequation}
\gamma_k^{(t)}(\mathbf{x}_i)=
\frac{\pi_k^{(t)}\,q(\mathbf{x}_i\mid\boldsymbol{\theta}_k^{(t)})}
{\sum_{j=1}^{K}\pi_j^{(t)}\,q(\mathbf{x}_i\mid\boldsymbol{\theta}_j^{(t)})}.
\end{myequation}
This step updates our beliefs about cluster membership: it computes the posterior
distribution of $z_i$ given $\mathbf{x}_i$ under the current model parameters. The
responsibilities $\gamma_k^{(t)}(\mathbf{x}_i)$ can be thought of as fractional labels: rather
than deciding definitively which component generated $\mathbf{x}_i$, the E step distributes
one unit of ``membership'' across all components in proportion to how well each component
explains the observation.

\subsubsection*{Updating parameters (the M step)}

The M step updates the parameters by maximising the expected complete-data
log-likelihood under the responsibilities just computed:
\begin{myequation}
\mathbb{E}_{\gamma}\!\left[\ell(\boldsymbol{\pi},\boldsymbol{\Theta}\mid\mathbf{X})\right]
=\sum_{i=1}^{n}\sum_{k=1}^{K}\gamma_k(\mathbf{x}_i)\log q(\mathbf{x}_i\mid\boldsymbol{\theta}_k)
+\sum_{k=1}^{K}\sum_{i=1}^{n}\gamma_k(\mathbf{x}_i)\log \pi_k.
\end{myequation}
The key point is that the responsibilities act as fractional membership indicators: each
observation contributes to each component, but with a weight equal to how strongly the
model associates that observation with that component. The M step thereby turns the
intractable incomplete-data maximisation into a weighted version of standard maximum
likelihood, which is typically tractable in closed form.

A fundamental property of the EM algorithm is monotonicity: each full iteration---one E
step followed by one M step---does not decrease the log-likelihood:
\begin{myequation}
\ell^{(t+1)}\geq \ell^{(t)}.
\end{myequation}
This guarantees convergence, though generally only to a \textbf{local} maximum of the
likelihood. In practice, convergence is often reached in a moderate number of iterations,
but different initializations can lead to different local optima. For this reason, it is common
to run EM multiple times from different starting points and retain the solution with the
largest achieved likelihood.

\subsection*{Gaussian mixture models (GMM)}

\subsubsection*{Model definition}

The most widely used mixture model for clustering in $\mathbb{R}^d$ is the \alert{Gaussian
mixture model} (GMM), in which each component is a multivariate Gaussian:
\begin{myequation}
q(\mathbf{x}\mid\boldsymbol{\theta}_k)=\mathcal{N}(\mathbf{x}\mid\boldsymbol{\mu}_k,\boldsymbol{\Sigma}_k)
=\frac{1}{(2\pi)^{d/2}|\boldsymbol{\Sigma}_k|^{1/2}}
\exp\!\left(-\frac{1}{2}(\mathbf{x}-\boldsymbol{\mu}_k)^T\boldsymbol{\Sigma}_k^{-1}(\mathbf{x}-\boldsymbol{\mu}_k)\right),
\end{myequation}
with parameters $\boldsymbol{\theta}_k=(\boldsymbol{\mu}_k,\boldsymbol{\Sigma}_k)$, a
mean vector and a covariance matrix. Compared with $k$-means, the GMM is a strictly
richer model: the covariance matrix $\boldsymbol{\Sigma}_k$ allows each component to be
elliptical, oriented, and of a size different from the others. In $k$-means, the implicit
assumption of spherical clusters of comparable variance corresponds to taking all
$\boldsymbol{\Sigma}_k = \sigma^2 \mathbf{I}$, so the GMM with general covariances is a
strict generalisation.

\subsubsection*{EM for GMMs}

With the Gaussian choice, the E step takes the form
\begin{myequation}
\gamma_k^{(t)}(\mathbf{x}_i)=
\frac{\pi_k^{(t)}\,\mathcal{N}(\mathbf{x}_i\mid\boldsymbol{\mu}_k^{(t)},\boldsymbol{\Sigma}_k^{(t)})}
{\sum_{j=1}^{K}\pi_j^{(t)}\,\mathcal{N}(\mathbf{x}_i\mid\boldsymbol{\mu}_j^{(t)},\boldsymbol{\Sigma}_j^{(t)})}.
\end{myequation}
To express the M step compactly, we define the effective (soft) cluster size
\begin{myequation}
n_k=\sum_{i=1}^{n}\gamma_k^{(t)}(\mathbf{x}_i).
\end{myequation}
Note that $n_k$ is in general not an integer: it represents the total fractional membership
allocated to component $k$ across all data points. The M-step updates are then weighted
versions of the usual sample estimators:
\begin{myequation}
\pi_k^{(t+1)}=\frac{1}{n}\sum_{i=1}^{n}\gamma_k^{(t)}(\mathbf{x}_i)=\frac{n_k}{n},
\end{myequation}
\begin{myequation}
\boldsymbol{\mu}_k^{(t+1)}=
\frac{\sum_{i=1}^{n}\gamma_k^{(t)}(\mathbf{x}_i)\mathbf{x}_i}{\sum_{i=1}^{n}\gamma_k^{(t)}(\mathbf{x}_i)}
=\frac{1}{n_k}\sum_{i=1}^{n}\gamma_k^{(t)}(\mathbf{x}_i)\mathbf{x}_i,
\end{myequation}
\begin{myequation}
\boldsymbol{\Sigma}_k^{(t+1)}=
\frac{1}{n_k}\sum_{i=1}^{n}\gamma_k^{(t)}(\mathbf{x}_i)\,
(\mathbf{x}_i-\boldsymbol{\mu}_k^{(t+1)})(\mathbf{x}_i-\boldsymbol{\mu}_k^{(t+1)})^T.
\end{myequation}
These formulas make the role of responsibilities transparent: points with large $\gamma_k$
strongly influence the parameters of component $k$, while points with small $\gamma_k$
contribute only marginally. The structure of the updates is therefore a natural weighted
generalisation of the standard sample mean and sample covariance, with the
responsibilities playing the role of the observation weights.

\subsection*{Poisson mixture models}

The mixture framework is not limited to continuous data. For \emph{count} observations, a
Poisson distribution is a natural choice for the component density. In the simplest
univariate case, the component distribution is
\begin{myequation}
q(x\mid\boldsymbol{\theta}_k)=q(x\mid\lambda_k)=\frac{e^{-\lambda_k}\lambda_k^x}{x!},
\end{myequation}
where $\lambda_k>0$ is the rate parameter of component $k$. The EM algorithm for
Poisson mixtures mirrors the general case. The E step is identical in structure:
\begin{myequation}
\gamma_k^{(t)}(x_i)=
\frac{\pi_k^{(t)}\,q(x_i\mid\lambda_k^{(t)})}{\sum_{j=1}^{K}\pi_j^{(t)}\,q(x_i\mid\lambda_j^{(t)})}.
\end{myequation}
In the M step, the mixing coefficients update exactly as before, giving $\pi_k^{(t+1)}=
n_k/n$, and maximising the expected log-likelihood with respect to $\lambda_k$ yields
\begin{myequation}
\lambda_k^{(t+1)}=
\frac{\sum_{i=1}^{n}\gamma_k^{(t)}(x_i)\,x_i}{\sum_{i=1}^{n}\gamma_k^{(t)}(x_i)}
=\frac{1}{n_k}\sum_{i=1}^{n}\gamma_k^{(t)}(x_i)\,x_i.
\end{myequation}
The updated rate is therefore the \emph{responsibility-weighted} sample mean of the
counts: an intuitively natural result, since the sample mean is the maximum likelihood
estimator for the Poisson rate in the fully observed case. The Poisson mixture illustrates
how the EM algorithm extends gracefully to non-Gaussian component families, with the
same structural logic and the same closed-form convenience.

\subsection*{Model selection for mixture models}

Once mixture models have been fitted for different values of $K$, we still need to decide
how many components provide an adequate description of the data. Because mixture
models define a likelihood, we can use information criteria that explicitly trade off fit
quality against model complexity. The \textbf{Akaike information criterion (AIC)} and the
\textbf{Bayesian information criterion (BIC)} are:
\begin{myequation}
\text{AIC}(K)=2\cdot(\text{number of parameters})-2\ln L,
\end{myequation}
\begin{myequation}
\text{BIC}(K)=(\text{number of parameters})\cdot\ln n-2\ln L.
\end{myequation}
Both criteria penalise models with more parameters, but BIC imposes a stronger penalty
when $n$ is large, reflecting the Bayesian argument that a model with more parameters
requires more evidence from the data to justify. For a GMM with $K$ components in $d$
dimensions, the number of free parameters is $K\big(d+d(d+1)/2\big)+K-1$, accounting
for $K$ mean vectors, $K$ covariance matrices (each with $d(d+1)/2$ free entries), and
$K$ mixing coefficients subject to one linear constraint. Lower values of AIC or BIC
indicate a preferable model according to the respective criterion.

The key advantage of these criteria over the heuristics used for $k$-means is that they are
\emph{derived} from a coherent statistical model rather than from geometric intuition.
Minimising AIC or BIC provides a principled answer to the question of how many
components are needed, and BIC in particular is asymptotically consistent: under suitable
regularity conditions and correct model specification, it selects the true number of
components as $n\to\infty$.

\subsection*{GMM vs.\ $k$-means: key differences}

Although both $k$-means and Gaussian mixture models are widely used to discover
structure in unlabeled data, they embody two rather different views of what a cluster is and
how it should be inferred. Understanding these differences is important both for choosing
between the two methods and for interpreting their respective outputs.

In $k$-means, a cluster is essentially a geometric region of the feature space: each
observation is assigned to exactly one prototype, and the outcome is a \emph{hard}
partition of the dataset. Once the algorithm has converged, every point belongs to one and
only one cluster, with no explicit representation of ambiguity. A GMM, by contrast, treats
cluster membership as a latent random variable. Rather than producing a binary decision, it
returns \emph{soft} assignments in the form of responsibilities $\gamma_k(\mathbf{x})$,
which can be interpreted as posterior probabilities. This probabilistic output is particularly
informative when clusters overlap or when points lie near boundaries, because it provides
an explicit quantification of uncertainty that $k$-means lacks entirely.

The difference in perspective is also reflected in the learning objectives. The $k$-means
algorithm minimises a purely geometric criterion, the within-cluster sum of squares $J$,
which encourages clusters to be compact around their centroids. GMMs are fitted by
maximising the likelihood $L$ under a generative model, thereby shifting from a
distance-based cost to a likelihood-based criterion. This shift has immediate consequences
for model selection: while $k$-means relies on heuristic procedures such as the elbow
method or the silhouette coefficient, GMMs can be compared in a principled way using
information criteria such as AIC and BIC, which explicitly balance goodness of fit against
model complexity.

Another important distinction concerns the shapes of clusters that each method can
represent. By construction, $k$-means tends to produce roughly spherical clusters of
comparable scale, since it relies on Euclidean distances to centroids. GMMs are more
flexible: each component has its own covariance matrix $\boldsymbol{\Sigma}_k$, allowing
clusters to be elliptical, oriented, and of different sizes and densities. This added flexibility
comes at a computational cost: the $k$-means updates are fast and typically require only a
few iterations, while fitting a GMM involves running EM with more expensive iterations that
may converge more slowly, especially in high dimensions or with full covariances.

Finally, the two methods differ in interpretability. A $k$-means solution is best understood
as a geometric partition induced by distances to prototypes, whereas a GMM provides a
fully generative probabilistic description of the data: clusters are not just regions of space
but components of a mixture distribution, and assignments are naturally probabilistic. This
makes GMMs particularly appealing when uncertainty matters, when cluster shapes are
not spherical, or when one wants likelihood-based tools for comparing models with different
values of $K$.

\subsection*{Practical considerations}

Mixture models are especially appealing when one needs more than a partition: they
provide uncertainty-aware membership, a coherent probabilistic interpretation, and
likelihood-based tools for model comparison. They also naturally accommodate clusters
with different volumes and orientations (in the Gaussian case) and can be adapted to
non-Gaussian data by choosing appropriate component families.

These advantages come with genuine costs. EM is more computationally demanding than
$k$-means, and its convergence to a local optimum makes initialisation important; multiple
random starts are often necessary. In GMMs with full covariances, one must also be aware
of degeneracies such as near-singular covariance matrices, which can arise when a
component collapses around very few points. More generally, richer models carry more
parameters, so reliable estimation requires sufficient data, and the chosen component
family must be a reasonable approximation to the true data-generating process.

A common practical workflow reflects this trade-off: begin with $k$-means as a fast
exploratory tool and as a source of good initializations for EM; move to mixture models
when soft assignments or a likelihood-based interpretation are needed; standardise features
before fitting Gaussian components; and compare several values of $K$ through AIC and
BIC, complementing these criteria with domain knowledge and qualitative inspection of the
resulting clusters.
%\end{document}